\documentclass[11pt,a4paper]{article}

\usepackage[a4paper,top=1.35cm,bottom=1.45cm,left=1.55cm,right=1.55cm]{geometry}
\usepackage{fontspec}
\usepackage{polyglossia}
\setmainlanguage{russian}
\setmainfont{Times New Roman}[Ligatures=TeX]
\newfontfamily\cyrillicfont{Times New Roman}
\setsansfont{Arial}[Ligatures=TeX]
\newfontfamily\cyrillicfontsf{Arial}
\setmonofont{Consolas}
\newfontfamily\cyrillicfonttt{Consolas}

\usepackage{amsmath}
\usepackage{amssymb}
\usepackage{booktabs}
\usepackage{array}
\usepackage{tabularx}
\usepackage{graphicx}
\usepackage{xcolor}

\definecolor{msublue}{RGB}{24,58,104}
\definecolor{msured}{RGB}{143,32,46}
\definecolor{softblue}{RGB}{238,243,249}
\definecolor{darkgray}{RGB}{45,48,53}

\setlength{\parindent}{1.1cm}
\setlength{\parskip}{0pt}
\setlength{\abovedisplayskip}{5pt}
\setlength{\belowdisplayskip}{5pt}
\renewcommand{\arraystretch}{1.12}
\color{darkgray}

\newcommand{\doctitle}[1]{%
  {\centering
    {\fontsize{18}{21}\selectfont\bfseries\color{msublue} #1\par}
    \vspace{0.18cm}
    {\large Краткий отчёт о проделанной работе\par}
    \vspace{0.12cm}
    {\normalsize Подъяблонский Егор Петрович, ВМК МГУ, 2026\par}
    \vspace{0.25cm}
    \color{msured}\rule{\textwidth}{0.8pt}\par
  }
}

\newcommand{\sectiontitle}[1]{%
  \vspace{0.18cm}
  {\large\bfseries\color{msublue} #1}\par\vspace{0.08cm}
}

\newcommand{\subsectiontitle}[1]{%
  \vspace{0.10cm}
  {\normalsize\bfseries\color{msured} #1}\par\vspace{0.04cm}
}

\newcommand{\note}[1]{%
  \begingroup
  \setlength{\fboxsep}{6pt}
  \noindent\colorbox{softblue}{\parbox{\dimexpr\textwidth-2\fboxsep\relax}{#1}}
  \endgroup
}

\begin{document}
\pagestyle{plain}

\doctitle{Квантизация диффузионных моделей машинного обучения}

\sectiontitle{1. Задача исследования}

Диффузионные модели обеспечивают высокое качество генерации изображений, но требуют значительного объёма памяти для хранения весов и больших вычислительных ресурсов. Цель работы состоит в исследовании post-training quantization для Stable Diffusion 1.5 и построении векторного метода, который сохраняет качество сильного VQDM, но избегает его наиболее дорогой части \textemdash{} совместного поиска комбинаций кодовых слов.

Основной объект экспериментов \textemdash{} U-Net модели Stable Diffusion 1.5. На текущем этапе квантуется \texttt{mid\_block}: полный набор из 18 Linear/Conv2d-слоёв, а также VQDM-style набор из 16 слоёв, в котором исключены две проекции \texttt{time\_emb\_proj}. Квантуются только веса; последующий fine-tuning модели не применяется.

\sectiontitle{2. Экспериментальный протокол}

Для всех методов используются одинаковые prompts, seed, параметры сэмплирования и набор из 96 изображений на один вариант. Скалярные методы сравниваются при 2, 3 и 4 битах на вес. Для векторных методов задаются целевые профили 2.15, 3.15 и 4.15 bpp; фактическая битность пересчитывается с учётом индексов и хранения кодбуков.

Калибровочные активации собираются на фиксированном наборе текстовых запросов и шагов диффузии. Они используются для оценки того, насколько ошибка в конкретном весе влияет на выход слоя. Благодаря этому сравниваются не только приближения весовой матрицы, но и их поведение на характерных входах модели.

\sectiontitle{3. Метрики}

\begin{itemize}\setlength\itemsep{0.12cm}\setlength\topsep{0.06cm}
  \item \textbf{CLIP I-I} $\uparrow$ \textemdash{} семантическое сходство изображений исходной и квантизованной моделей.
  \item \textbf{LPIPS} $\downarrow$ \textemdash{} перцептивное расстояние между теми же изображениями.
  \item \textbf{CLIP T-I} $\uparrow$ \textemdash{} соответствие сгенерированного изображения исходному текстовому запросу.
  \item \textbf{PSNR и SSIM} $\uparrow$ \textemdash{} пиксельное и структурное сходство с результатом исходной модели.
  \item \textbf{bpp} \textemdash{} фактическое число бит на один квантуемый вес; \textbf{$T_q$} \textemdash{} время построения квантизованного представления с необходимой калибровкой.
\end{itemize}

Для компактного ранжирования используется
\[
  \mathrm{Score}=\mathrm{CLIP}_{I-I}-\mathrm{LPIPS}.
\]
Эта величина одновременно поощряет семантическое сходство и штрафует перцептивное расхождение, но окончательные выводы проверяются по всем метрикам, а не только по Score.

\note{Все приведённые ниже результаты относятся к единому no-fine-tuning протоколу на mid-block. Это позволяет сравнивать алгоритмы в одинаковых условиях, но пока не является оценкой квантизации полного U-Net.}

\newpage

\sectiontitle{4. Рассмотренные базовые методы}

\subsectiontitle{Скалярная квантизация}

\begin{itemize}\setlength\itemsep{0.28cm}\setlength\topsep{0.06cm}
  \item \textbf{Scalar PTQ.} Каждый вес независимо округляется к ближайшему уровню равномерной сетки. Метод практически не требует обучения и служит нижней контрольной точкой, но не учитывает распределение активаций и структуру диффузионной модели.

  \item \textbf{Q-Diffusion-like.} Диапазоны и параметры скалярной сетки выбираются на калибровочных данных с учётом ошибки восстановления выхода блока. Это делает квантизацию чувствительной к реальному прохождению сигнала через диффузионную модель.

  \item \textbf{PTQ4DM-like.} Используются timestep-aware калибровочные данные, послойная или поблочная политика настройки и reconstruction objective. В отличие от обычной PTQ, калибровка учитывает разные уровни шума в процессе денойзинга.

  \item \textbf{BitsFusion-lite.} Слоям назначается разная скалярная битность при заданном среднем bpp: чувствительные слои получают больше бит, устойчивые \textemdash{} меньше. В эксперименте используется calibration-only версия без обучаемой distillation-части исходного метода.
\end{itemize}

\subsectiontitle{Векторная квантизация}

\begin{itemize}\setlength\itemsep{0.28cm}\setlength\topsep{0.06cm}
  \item \textbf{Vanilla VQ.} Короткий вектор весов заменяется ближайшим вектором одного кодбука. Метод очень быстро строится, но единственный индекс ограничивает выразительность аппроксимации.

  \item \textbf{AQUATIC-Diff-like.} Используется аддитивное представление несколькими кодбуками и калибровка по активациям. Реализованная версия проверяет основную codebook-идею в общем протоколе, но не включает все обучаемые процедуры полной AQUATIC-Diff.

  \item \textbf{VQDM.} Весовой вектор представляется суммой кодовых слов из нескольких кодбуков. Индексы выбираются совместно с помощью beam search, а качество комбинации оценивается activation-aware функционалом. VQDM даёт наиболее точную аппроксимацию среди базовых методов, однако стоимость поиска быстро растёт с числом кодбуков и квантуемых слоёв.
\end{itemize}

\sectiontitle{5. Интерпретация набора бейзлайнов}

Scalar PTQ показывает нижнюю границу сложности и качества. Q-Diffusion-like и PTQ4DM-like оценивают пользу diffusion-aware калибровки. BitsFusion-lite проверяет неравномерное распределение скалярного битового бюджета. Vanilla VQ является нижним векторным бейзлайном, AQUATIC-Diff-like \textemdash{} более содержательным аддитивным вариантом, а VQDM задаёт верхнюю точку качества для activation-aware многокодбуковой квантизации.

\note{Суффиксы \textit{like/lite} означают адаптацию основной идеи статьи к SD1.5 и единому no-fine-tuning протоколу. Результаты таких вариантов не следует трактовать как полную репродукцию обучаемых версий исходных методов.}

\newpage

\sectiontitle{6. Предлагаемый activation-aware RVQ}

Пусть фрагмент весов представлен вектором $w\in\mathbb{R}^{d}$. В residual vector quantization он аппроксимируется суммой кодовых слов из последовательности кодбуков:
\[
  \hat w=\sum_{m=1}^{M}C_m[z_m],
  \qquad r_0=w,\qquad r_m=r_{m-1}-C_m[z_m].
\]
Каждая стадия квантует остаток предыдущей. В VQDM комбинация $(z_1,\ldots,z_M)$ выбирается совместно с beam search; в RVQ индексы находятся последовательно. Такая жадность может ухудшать качество, зато вычислительная стоимость растёт существенно медленнее и позволяет исследовать больше кодбуков и профилей.

\subsectiontitle{Activation-aware функционал}

Для слоя с весами $W$ и матрицей калибровочных входов $X$ минимизируется ошибка выхода слоя:
\[
  \mathcal L(\hat W)=\|X(W-\hat W)^\top\|_F^2
  =\operatorname{tr}\!\left((W-\hat W)X^\top X(W-\hat W)^\top\right).
\]
В отличие от weight-only VQ, здесь одинаковые по норме ошибки весов могут иметь разную стоимость: матрица $X^\top X$ усиливает направления, которые чаще и сильнее используются активациями.

\subsectiontitle{Распределение размера кодбуков по стадиям}

Вместо одинакового $K$ для всех стадий используются значения $K_1,\ldots,K_M$. Битовая стоимость индексов равна
\[
  b_{\mathrm{index}}=\frac{1}{d}\sum_{m=1}^{M}\log_2 K_m.
\]
Поэтому ёмкость можно перераспределять между ранними и поздними стадиями почти без изменения общего бюджета. Это отличается от простого увеличения всех кодбуков, которое автоматически повышает bpp.

\subsectiontitle{Ортогональное преобразование Hadamard}

Перед RVQ координаты блока могут быть перемешаны ортогональной матрицей Hadamard. Норма вектора при этом сохраняется, но меняется геометрия дискретной задачи: энергия и корреляции распределяются между координатами, что иногда облегчает работу коротких кодовых слов. Преобразование не включается безусловно, поскольку его эффект зависит от битового профиля и размера блока.

\subsectiontitle{Post-RVQ recoding по полной матрице $X^\top X$}

При обучении кодбуков activation-aware функционал обрабатывается по коротким блокам. После этого кодбуки фиксируются, а дискретные индексы дополнительно уточняются coordinate descent с использованием полной матрицы $X^\top X$. Это позволяет учитывать взаимодействия между блоками. Такой шаг не является fine-tuning: веса модели и векторы кодбуков не обучаются градиентно, изменяются только индексы уже построенного представления.

\note{Итоговая рабочая версия: $d=8$, разные $K_m$ по стадиям, full-$X^\top X$ recoding и отключённый local recode. Hadamard выбирается отдельно для каждого bpp-профиля.}

\newpage

\sectiontitle{7. Промежуточные результаты}

Ниже приведено сравнение на полном наборе из 18 слоёв mid-block. Для каждого битового профиля показаны VQDM, предлагаемый RVQ-AW и лучший следующий бейзлайн по Score.

\begin{center}
\resizebox{\textwidth}{!}{%
\begin{tabular}{c l r r r r r r}
  \toprule
  Профиль & Метод & bpp & CLIP I-I & CLIP T-I & LPIPS & Score & $T_q$, с \\
  \midrule
  2 & VQDM       & 2.02 & 0.831 & 0.305 & 0.244 & \textbf{0.586} & 1857.1 \\
  2 & \textbf{RVQ-AW} & 2.03 & 0.819 & 0.305 & 0.254 & 0.564 & 604.2 \\
  2 & BitsFusion & 2.05 & 0.768 & 0.312 & 0.300 & 0.468 & 88.4 \\
  \midrule
  3 & VQDM       & 3.02 & 0.885 & 0.308 & 0.164 & \textbf{0.721} & 2846.8 \\
  3 & \textbf{RVQ-AW} & 3.02 & 0.879 & 0.308 & 0.175 & 0.704 & 992.8 \\
  3 & Vanilla VQ & 3.01 & 0.828 & 0.319 & 0.212 & 0.617 & 1.91 \\
  \midrule
  4 & VQDM       & 4.03 & 0.926 & 0.310 & 0.101 & \textbf{0.825} & 3314.0 \\
  4 & \textbf{RVQ-AW} & 4.03 & 0.911 & 0.311 & 0.125 & 0.786 & 1400.0 \\
  4 & Vanilla VQ & 4.00 & 0.881 & 0.317 & 0.144 & 0.737 & 1.57 \\
  \bottomrule
\end{tabular}}
\end{center}

\subsectiontitle{Результаты проверок отдельных гипотез}

\begin{center}
\small
\begin{tabularx}{\textwidth}{>{\bfseries}p{3.8cm}X}
  \toprule
  Проверка & Наблюдение \\
  \midrule
  Hadamard & H16 улучшил Score на 0.014 в чистом эксперименте для 3.15 bpp, но при 2.15 и 4.15 лучшим остался вариант без преобразования. \\
  Full-$X^\top X$ recoding & Внутренний функционал уменьшился в среднем на 47.53\%; влияние на внешние image-метрики оказалось положительным в среднем, но не для каждой конфигурации. \\
  Размерность $d$ & При $K=256$ вариант $d=16$ проиграл $d=8$ по качеству и потребовал почти вдвое больше времени. \\
  Local recode & Уточнение после промежуточных стадий выиграло только 2 из 20 парных сравнений; среднее изменение Score составило $-0.0408$. \\
  \bottomrule
\end{tabularx}
\end{center}

\subsectiontitle{Итоговые профили RVQ-AW}

\begin{center}
\begin{tabular}{c c c c c}
  \toprule
  Цель & Hadamard & $(K_1,\ldots,K_M)$ & $d$ & Фактический bpp \\
  \midrule
  2.15 & noH & $[64,1024]$ & 8 & 2.03 \\
  3.15 & H32 & $[128,256,512]$ & 8 & 3.02 \\
  4.15 & H16 & $[512,512,128,128]$ & 8 & 4.03 \\
  \bottomrule
\end{tabular}
\end{center}

\sectiontitle{8. Основной вывод и дальнейшая работа}

RVQ-AW занимает второе место во всех трёх профилях. Разрыв с VQDM по Score составляет 0.022, 0.017 и 0.039 для профилей около 2, 3 и 4 bpp соответственно. При этом RVQ-AW превосходит остальные проверенные методы; например, около 3 bpp он достигает Score 0.704 против 0.617 у следующего бейзлайна. Таким образом, последовательная остаточная квантизация сохраняет большую часть преимущества VQDM без совместного AVQ/beam-search поиска.

Следующий этап исследования \textemdash{} контролируемый совместный перебор $K$, $d$, Hadamard и вариантов recoding, timestep-aware калибровка, переход от mid-block к полному U-Net и оценка на более сильных генеративных метриках. Fine-tuning или distillation следует рассматривать отдельно, чтобы не смешивать post-training и обучаемые методы.

\end{document}
